% =====================================================================
%  chunks/ch03h.tex —— 原书 p.103–112（PDF 110–119）
%  第 3 章 §7 经典同伦论机器。Leray 谱序列。纤维丛的同调论。
%  Cartan--Serre 方法。Postnikov 塔。Adams 谱序列
% =====================================================================
% [原文 p.103]
\begin{enumerate}[label=(\roman*),leftmargin=2.2em,start=2]
\item $\Z/p$-代数 $A$ 是自由左 $C$-模，即存在由元素 $a_i\in A^{n_i}$（$i=0,1,2,\dots$，$0=n_0<n_1\le n_2\le n_3\le\cdots$，$a_0=1$）组成的基，使 $A$ 的每个元素 $a$ 都可唯一地写成
\[
a=\sum_{i\ge0}c_ia_i,\qquad c_i\in C.
\]
\end{enumerate}

在上述条件下，存在一个谱序列（它是 20 世纪 50 年代中所发现的 Serre--Hochschild 谱序列的特例）：它是「纤维丛」$f:A\to B$ 的、以 $C$ 为「纤维」的谱序列，模仿 CW 复形纤维丛的 Leray 谱序列的某些性质（见 §7）。它具有三阶分次的代数谱序列的形式
\[
E_m^{q,l,r},\qquad d_m:E_m^{q,l,r}\longrightarrow E_m^{q+m,l-m+1,r},\qquad d_m^2=0,
\]
其中：
\begin{enumerate}[label=\arabic*.,leftmargin=2em]
\item $E_{m+1}^{*,*,*}=H^*(E_m^{*,*,*},d_m)$，$d_m(uv)=d_m(u)v\pm ud_m(v)$；
\item $E_2^{q,l,r}=\sum\limits_{r_1+r_2=r}H^{q,r_1}(B;H^{l,r_2}(C))=\sum\limits_{r_1+r_2=r}H^{q,r_1}(B)\otimes_{\Z/p}H^{l,r_2}(C)$；
\item $\bigoplus\limits_{q+l=n}E_\infty^{q,l,r}\cong GH^{n,r}(A)$，
\end{enumerate}
其中 $GH^{n,r}(A)$ 表示关于产生上述谱序列的某个滤过的相伴环。指标 $r_1,r_2,r_3$ 所标出的分次分别对应于 $A$、$B$、$C$ 的分次。

上面提到的 Steenrod 运算 $St^i_p$（$p=2$ 时为 Steenrod 平方 $Sq^i$）的类似物（见 (6.14)）与上述谱序列的关系，与通常的 Steenrod 运算同纤维丛的 Leray 谱序列的关系大致相同；后者我们将在下一节考虑。

\section[经典同伦论机器。 Leray 谱序列。 Cartan--Serre 方法]{经典同伦论机器。Leray 谱序列。纤维丛的同调论。Cartan--Serre 方法。Postnikov 塔。Adams 谱序列}

\takeaway{本节把同伦论的四件经典机器串成一条线：Leray 谱序列把纤维丛的上同调拆成底与纤维的两族数据；Cartan--Serre 方法用 Hurewicz 定理与纤维化逐维算出同伦群；Postnikov 塔把这些群连同 $k$-不变量组织成空间本身；Adams 谱序列则在稳定范畴中把同伦群的计算化为 Steenrod 模的同调代数。}

「谱序列」方法由 Leray 于 20 世纪 40 年代中在关于空间映射的问题中首先发现，因而特别适用于纤维丛；作为同调代数中的有力工具，它具有根本的重要性，使人们（除其它用途之外）能够通过避免对几何结构作直接的细致考察，来计算一大类空间的同调群。

% [原文 p.104]
设 $p:Y\to X$ 是拓扑空间之间的映射。对每个 $j\ge0$，映射 $p$ 在 $Y$ 上确定一个预层 $\mathcal{F}^j$，它由
\begin{equation}
\mathcal{F}^j_U=H^j(p^{-1}(U);D),\qquad j\ge0, \tag{7.1}
\end{equation}
给出，其中 $U$ 取遍 $X$ 的开集，$D$ 是任一固定的 Abel 群。于是又得到上同调群
\[
E_2^{q,l}=H^q(X;\mathcal{F}^l)
\]
（当 $q<0$ 或 $l<0$ 时定义为零）。\emph{Leray 定理}断言：存在一个收敛到 $H^m(Y;D)$ 的谱序列
\begin{equation}
E_m^{q,l},\qquad d_m:E_m^{q,l}\longrightarrow E_m^{q+m,l-m+1},\qquad d_m^2=0, \tag{7.2}
\end{equation}
其中
\[
E_{m+1}^{q,l}=H^*(E_m^{q,l},d_m)\cong\Ker d_m/\Image d_m,
\]
\[
E_\infty^{q,l}=E_m^{q,l}\qquad\text{当 } m>q+l,
\]
使得
\begin{equation}
\sum_{q+l=m}E_\infty^{q,l}\cong GH^m(Y,D)=\sum_{l=0}^{m}B^{m-l}/B^{m-l-1}. \tag{7.3}
\end{equation}
这里 $GH^m(Y,D)$ 表示 $H^m(Y,D)$ 关于某个滤过的相伴群
\[
B^m=H^m(Y;D)\supset B^{m-1}\supset B^{m-2}\supset\cdots\supset B^0\supset0.
\]
对某些指标 $q,l$，群 $E_\infty^{q,l}$ 容易辨认；例如 $E_\infty^{0,m}=B^0\subset H^m(Y;D)$ 正是投影同态 $p^*$ 的像：
\begin{equation}
E_\infty^{0,m}=B^0=\Image p^*,\qquad p^*:H^m(X;D)\longrightarrow H^m(Y;D). \tag{7.4}
\end{equation}

% [Tutor Note: (7.1) 中的预层一般不是层，因而 $H^q(X;\mathcal{F}^l)$ 须按层上同调（即把预层层化后再算）理解；这正是 $E_2$ 项在 $\pi_1(X)$ 非平凡时要以局部系数计算的原因。]

对以 $F$ 为纤维的局部平凡纤维化 $p:Y\to X$，上述 \emph{Leray 谱序列}可以描述如下。设底空间 $X$ 是带 $n$-骨架 $X^n$（$n=0,1,2,\dots$）的 CW 复形。$H^m(Y;D)$ 的滤过由空间 $Y$ 的下述自然滤过决定：
\[
Y_n=p^{-1}(X^n),\qquad Y_n\xrightarrow{\phi_n}Y,
\]
\[
Y_0\subset Y_1\subset Y_2\subset\cdots\subset Y_n\subset\cdots,\qquad Y=Y_\infty.
\]
由包含映射 $\phi_k:Y_k\to Y$，对每个 $m$ 得到诱导同态 $\phi_k^*:H^m(Y)\to H^m(Y_k)$；显然当 $k=m$ 时 $\phi_k^*$ 是同构。所求的滤过
\[
H^m(Y;D)=B^m\supset B^{m-1}\supset B^{m-2}\supset\cdots\supset B^0\supset0
\]

% [原文 p.105]
通过令
\[
B^k=\Ker\{\phi^*_{m-k}:H^m(Y;D)\to H^m(Y_{m-k};D)\},\qquad k=m,m-1,\dots,0
\]
得到。此外，若底空间 $X$ 连通且只有一个顶点（0-胞腔）——这总可以不失一般性地假定——则 (7.3) 中的商 $B^m/B^{m-1}$ 同构于 $H^m(Y;D)$ 在 $H^m(F;D)$ 中的像，后者由纤维 $F$ 到全空间 $Y$ 的包含诱导出同态：
\[
i:F\longrightarrow Y,\qquad F=p^{-1}(x_0),
\]
\begin{equation}
\Image\{i^*:H^m(Y;D)\to H^m(F;D)\}\cong B^m/B^{m-1}\subset H^m(F;D). \tag{7.5}
\end{equation}
进一步，对局部平凡的纤维丛（以及对 Serre 纤维化——事实上对任何具有覆叠同伦性质的纤维化），Leray 定理给出 $E_2^{q,l}$ 项的显式形式：
\begin{equation}
E_2^{q,l}=H^q(X;\{H^l(F;D),\rho\}). \tag{7.6}
\end{equation}
这里 $H^l(F;D)$ 上基本群 $\pi_1(X)$ 的表示 $\rho$ 由 $\pi_1(X)$ 在纤维 $F$ 上的作用决定。若 $\pi_1(X)$ 平凡地作用（特别地，若 $\pi_1(X)=0$），则由 (7.6) 可知 $E_2^{q,l}=H^q(X;H^l(F))$，即 $E_2^{q,l}$ 项与直积 $X\times F$ 的上同调一致：
\begin{equation}
\sum_{q,l}E_2^{q,l}\cong H^*(X\times F;D). \tag{7.7}
\end{equation}
对平凡丛 $Y=X\times F$（直积）有 $E_2^{q,l}=E_\infty^{q,l}$，即对一切 $m\ge2$ 微分 $d_m$ 都为零。

然而对非平凡纤维丛（或「斜积」）且底空间单连通的情形，由于存在非零微分 $d_m$（$m\ge2$），全空间的上同调变得与 $H^*(X\times F;D)$ 不同。于是对斜积而言，上同调比直积的「更少」。这并不奇怪，理由如下。假定纤维 $F$ 与底空间 $X$ 都是 CW 复形；由于 $X$ 的每个胞腔 $\sigma^m$ 上的丛是平凡的，全空间的胞腔分解与直积 $X\times F$ 的相同：
\[
p^{-1}(\sigma^m)\cong\sigma^m\times F.
\]
全空间 $Y$ 的上同调群 $H^*(Y;D)$ 中的「斜」是在应用边缘算子 $\partial$ 时出现的：
\[
\partial\sigma_Y^{q+l}=\partial(\sigma_X^q\times\sigma_F^l)=\sigma^q\times(\partial\sigma_F^l)+(\partial_1+\partial_2+\cdots)\sigma_Y^{q+l},
\]
其中 $\partial_j\sigma_Y^{q+l}$ 位于底空间 $X$ 的 $(q-j)$-骨架 $X^{q-j}$ 上方。算子 $\partial_1$ 可以用底空间的上同调（系数取自 $\pi_1(X)$-模 $H^*(F)$）显式算出（这是 Leray 的结果）。

% [Tutor Note: 读 (7.3)–(7.5) 时抓住「滤过—相伴群」这一对概念：$E_\infty$ 逐项给出的是商 $B^m/B^{m-1}$，而不是 $H^m$ 本身，故 $m$ 维上同调要由一条对角线 $q+l=m$ 上的全部 $E_\infty^{q,l}$ 拼回；若只知道 $E_\infty$ 项而不知道扩张问题，一般还不能唯一确定 $H^m$。]

% [原文 p.106]
若 $D$ 是环，则每个 $E_m^{*,*}$ 成为双分次环，且微分 $d_m$ 对双分次元素 $u,v$ 满足 $d_m(uv)=d_m(u)v\pm ud_m(v)$（所得环结构已含在 Leray 定理之中）。定理有完全类似的同调版本；在域上，上同调谱序列与同调谱序列互为对偶。

Leray 谱序列中的\emph{超渡同态} $\tau$ 如下定义：考虑对 $(Y,F)$ 与同态
\[
H^q(F)\xrightarrow{\delta}H^{q+1}(Y,F)\qquad(\text{上边缘同态}),
\]
\[
H^{q+1}(X,x_0)\xrightarrow{p^*}H^{q+1}(Y,F)\qquad(\text{投影同态}).
\]
复合 $(p^*)^{-1}\circ\delta=\tau$ 就是所求的超渡。注意 $\tau$ 不必在 $H^q(F)$ 上处处有定义，并且可能是多值的。若 $\tau(z)$ 有定义，则称 $H^q(F)$ 的元素 $z$ 是\emph{超渡的}。用纤维丛的 Leray 谱序列表述，超渡 $\tau$ 等价地定义为 $\tau=d_q$，其中
\begin{equation}
H^q(F)\supset E_q^{0,q}\xrightarrow{d_q}E_q^{q+1,0}=H^{q+1}(X)/\Ker p^*. \tag{7.8}
\end{equation}
由 (7.8) 可知，对超渡元素 $z\in H^q(F)$ 有 $d_j(z)=0$（$j<q$）。

当 $D=\Z/p$ 时，一切上同调环都是代数 $\mathcal{A}_p$ 上的模；结果表明，Leray 谱序列中的微分与 Steenrod 平方及幂交换；事实上，既然 $\delta$ 与 $p^*$ 与 Steenrod 平方及幂交换，微分 $d_q$ 也就与它们交换。这在某些困难的计算中具有关键的重要性。特别地，若 $z\in H^q(F)$ 是超渡的，则 $St^i_p(z)$ 也是超渡的，并且
\[
\tau St^i_p(z)=St^i_p\tau(z).
\]

\begin{example}{1}
\emph{Borel 定理}（20 世纪 50 年代初）：若 $H^*(F;\Q)$ 是 $\Q$ 上的外代数，且 $H^*(Y;\Q)=0$，则 $H^*(X;\Q)$ 是某些外生成元在超渡下的像所生成的多项式代数。在域 $\Z/2$ 上，把 $H^*(F;\Z/2)$ 为外代数这一假设换成如下要求：存在一列有限的超渡生成元 $v_j\in H^{n_j}(F;\Z/2)$，构成 $H^*(F;\Z/2)$ 的加法基：
\[
v_{j_1},v_{j_2},\dots,v_{j_k},\qquad j_1<j_2<\cdots<j_k.
\]
结论相同：环 $H^*(X;\Z/2)$ 是多项式环。借助这个定理，计算出了 $K(\Z/2,n)$ 型空间在 $\Z/2$ 上的上同调以及 Steenrod 代数 $\mathcal{A}_2$（Serre，20 世纪 50 年代初）。
\end{example}

\begin{example}{2}
利用 Leray 谱序列，可以毫无困难地算出下述上同调环：

% [原文 p.107]
\[
H^*(\RP^n;\Z/2);\qquad H^*(\RP^n;\Z);\qquad H^*(\CP^n;\Z);
\]
\[
H^*(S^\infty/\Z/m;\Z);\qquad H^*(S^\infty/\Z/m;\Z/m);
\]
\[
H^*(SO_n;\Q);\qquad H^*(U_n;\Z);\qquad H^*(Sp_n;\Z);
\]
\[
H^*(BSO_n;\Q);\qquad H^*(BU_n;\Z);\qquad H^*(BSO_n;\Z/2);
\]
\[
H^*(BO_n;\Z/2).
\]
这些计算借助下述纤维化进行：
\begin{enumerate}[label=\arabic*.,leftmargin=2em]
\item $S^{2n+1}\to\CP^n$（纤维 $S^1$）；
\item $\RP^{2n+1}\to\CP^n$（纤维 $S^1$）；
\item $S^{2n+1}/\Z/m\to\CP^n$（纤维 $S^1$）；
\item $SO_n\to S^{n-1}$（纤维 $SO_{n-1}$）；
\item $U_n\to S^{2n-1}$（纤维 $U_{n-1}$）；
\item $E\to BO_n$（纤维 $O_n$），$E\sim*$；
\item $E\to BU_n$（纤维 $U_n$），$E\sim*$。
\end{enumerate}
\end{example}

\begin{example}{3}
作为上述 Borel 定理的平凡特例，对空间 $X=K(\Z,n)$ 有
\[
H^*(X;\Q)=\begin{cases}\Lambda(u), & u\in H^n(X),\ \text{若 } n \text{ 为奇数},\\ \Q[u], & u\in H^n(X),\ \text{若 } n \text{ 为偶数}.\end{cases}
\]
对有限群 $\pi$，上同调环 $H^*(K(\pi,n);\Q)$ 为零。对有限生成的 Abel 群 $\pi$，上同调环 $H^*(K(\pi,n);R)$（$R$ 为任意系数环）是有限生成的。
\end{example}

第 1 章中我们描述过：如何借助一个保持同伦型的构造，把任意映射 $f:X\to Y$ 变成 Serre 纤维化
\[
p:\tilde X\longrightarrow\tilde Y,\qquad \tilde X\sim X,\quad \tilde Y\sim Y,\quad p\sim f,
\]
其中 $\tilde Y$ 是 $f$ 的映射柱，可缩到 $Y$；$\tilde X$ 是 $\tilde Y$ 中始于 $X$ 的点（把 $X$ 与 $X\times\{0\}\subset\tilde Y$ 等同）、终于 $Y$ 的点的道路空间；$p$ 把每条道路送到它在 $Y$ 中的终点。空间 $\tilde X$ 可缩到 $X$。

现在叙述计算单连通空间 $X$ 的同伦群 $\pi_j(X)$ 的 Cartan--Serre 方法。设当 $j<n$ 时 $\pi_j(X)=0$，并令 $Y=K(\pi_n(X),n)$。

% [原文 p.108]
于是存在映射 $f:X\to Y$，其诱导同态 $f_*:\pi_n(X)\to\pi_n(Y)$ 是同构。把上述构造用于映射 $f$，得到纤维丛（Serre 纤维化）
\[
p:\tilde X\longrightarrow\tilde Y\sim K(\pi_n(X),n).
\]
由此可知，纤维化 $p:\tilde X\to\tilde Y$ 的纤维 $F$ 在 $n$ 维以上的同伦群与 $X$ 的相同：
\begin{equation}
\begin{aligned}
\pi_j(F)&\cong\pi_{j+1}(X) &&\text{对 } j\ge n,\\
\pi_j(F)&=0 &&\text{对 } j<n.
\end{aligned}
\tag{7.9}
\end{equation}
因此算出 $\pi_*(F)$ 就等于算出 $\pi_*(X)$。Cartan--Serre 方法的最简单形式如下：

\para{第 1 步.} 用同调的 Leray 谱序列（施于纤维化 $p:\tilde X\to\tilde Y$）算出纤维 $F$ 的同调群 $H_*(F)$。

\para{第 2 步.} 用 Hurewicz 定理计算群 $\pi_{n+1}(X)$：
\[
H_n(F)\cong\pi_n(F)\cong\pi_{n+1}(X).
\]

\para{第 3 步.} 重复整个过程。

不过，计算单连通空间 $X$ 的同伦群还有更方便的程序。通过向空间 $X$ 粘上维数 $\ge q$ 的胞腔以「消去」维数 $\ge q$ 的同伦群，便得到空间 $X_q\supset X$，使
\begin{equation}
\begin{aligned}
\pi_j(X)&\cong\pi_j(X_q) &&\text{对 } j<q,\\
\pi_j(X_q)&=0 &&\text{对 } j\ge q.
\end{aligned}
\tag{7.10}
\end{equation}
因此若 $\pi_j(X)=0$（$j<n$），则对一切 $j$ 有 $\pi_j(X_n)=0$，即 $X_n$ 可缩；而 $X_{n+1}$ 是 $K(\pi,n)$ 型的，$\pi=\pi_n(X)$。考虑下述纤维化序列（由把相应的包含 $X_{i+1}\subset X_i$ 化为纤维化而得到）：
\begin{align*}
(n+1)&\qquad \tilde X_{n+2}\longrightarrow\tilde X_{n+1}\qquad(\text{纤维 } F_{n+1}\sim K(\pi_{n+1}(X),n+1));\\
(n+2)&\qquad \tilde X_{n+3}\longrightarrow\tilde X_{n+2}\qquad(\text{纤维 } F_{n+2}\sim K(\pi_{n+2}(X),n+2));\\
&\qquad\qquad\cdots\cdots\cdots\\
(n+k)&\qquad \tilde X_{n+k+1}\longrightarrow\tilde X_{n+k}\qquad(\text{纤维 } F_{n+k}\sim K(\pi_{n+k}(X),n+k)).
\end{align*}
现在，以 $K(\pi,m)$ 型空间为纤维、底空间 $B$ 单连通的纤维丛，在同伦意义下由单个上同调类

% [Tutor Note: 第 2 步之所以成立，是因为由 (7.9) 有 $\pi_n(F)=\pi_{n+1}(X)$，而 $F$ 是单连通的（$X$ 单连通、$n\ge2$），故 Hurewicz 定理给出 $H_n(F)\cong\pi_n(F)$；若 $n=1$，Hurewicz 只给出同调化，须另作处理。]

% [原文 p.109]
\begin{equation*}
u_m\in H^{m+1}(B,\pi)
\end{equation*}
决定。于是，上述每个纤维化 $\tilde X_{n+k+1}\to\tilde X_{n+k}$（纤维 $F_{n+k}\sim K(\pi_{n+k},n+k)$）都对应一个特定的上同调类
\[
u_{n+k}\in H^{n+k+1}(X_{n+k},\pi_{n+k}(X)).
\]
这些上同调类 $u_{n+k}$ 称为空间 $X$ 的 \emph{Postnikov 不变量}。以 $F_{n+k}\sim K(\pi_{n+k},n+k)$ 为纤维的纤维化序列
\begin{equation}
\cdots\longrightarrow\tilde X_{n+3}\longrightarrow\tilde X_{n+2}\longrightarrow\tilde X_{n+1}\longrightarrow* \tag{7.11}
\end{equation}
（这里 $*$ 表示可缩空间）称为 $X$ 的 \emph{Postnikov 塔}。由此可知，群组 $\pi_j(X)$ 与 Postnikov 不变量 $u_j$ 一起决定了单连通空间 $X$ 的同伦型。

Cartan--Serre 方法使人们能在递归地构造空间 $X_{n+j}$ 的过程中算出 Postnikov 不变量，其中要用到 Eilenberg--MacLane 空间 $K(\pi,m)$ 的上同调。对有理同伦群 $\pi_i(X)\otimes\Q$，由于上同调环 $H^*(K(\pi,m);\Q)$ 的结构简单，这一计算程序变得非常简单。

\begin{example*}{1}
对球面 $S^n$，群 $\pi_i(S^n)\otimes\Q$ 由
\[
\pi_n(S^n)\cong\Z,\qquad \pi_{4n-1}\otimes\Q=\Q
\]
给出，其余有理同伦群都平凡（由此可知在其它一切情形下 $\pi_i(S^n)$ 都是有限群）。
\end{example*}

\begin{example*}{2}
对单连通的有限 CW 复形 $X$，同伦群 $\pi_j(X)$ 都是有限生成的。（这两个结果都是 Serre 在 20 世纪 50 年代初得到的。）
\end{example*}

\begin{example*}{3}
\emph{Cartan--Serre 定理}（20 世纪 50 年代初）。若单连通空间 $X$ 的代数 $H^*(X;\Q)$ 是自由斜对称代数，则 $\pi_j(X)\otimes\Q$ 同构于 $H^j(X;\Q)$ 上那些在 $H_j(X;\Q)$ 的非平凡可分解元素上取零的线性 $\Q$-形式所成的向量空间。这个定理适用于 $H$-空间，因为按 Hopf 定理，这类空间的 $H^*(X;\Q)$ 是自由斜对称的。于是该结果特别适用于环路空间 $\Omega X$：对任意有限单连通 CW 复形 $X$，知道了 $H^*(\Omega X;\Q)$ 就完全知道了群 $\pi_j(X)\otimes\Q$。然而这样的上同调代数并不总能容易地具体算出。（$H$-空间的定义及其它细节见第 4 章 §2。）
\end{example*}

更晚些时候，在 20 世纪 70 年代初，Sullivan 把有理同伦型（「$\Q$-型」）的一般理论精确化，该理论适用于有限单连通 CW 复形，其中一切计算都在有理范畴中进行，即把一切不变量与 $\Q$ 作张量积。他证明了：这样的 CW 复形的有理同伦型由下述等价类决定，即某个微分斜对称代数的等价类。

% [原文 p.110]
通过 Kähler 几何，这特别导致如下结果：单连通「Kähler 流形」的 $\Q$-型完全由它的有理上同调代数决定（Deligne、Sullivan、Griffiths、Morgan，20 世纪 70 年代中）。

在某些特殊情形下，Cartan--Serre 方法允许精确计算同伦群，而不只是算到与 $\Q$ 作张量积为止。Adams 把 Cartan--Serre 方法重新组织成一个谱序列，即所谓「Adams 谱序列」。他注意到，对稳定同伦的计算，实际上使用所谓「Adams 分解」比使用 Postnikov 塔更方便。

设 $X$ 是维数小于 $2n-1$ 的有限 CW 复形（或同伦等价于这样的复形），且当 $j\le n-1$ 时 $\pi_j(X)=0$。现在把上同调环 $H^*(X;\Z/p)$ 看作 Steenrod 代数 $\mathcal{A}_p$ 上的模。设 $\{m_1,\dots,m_k\}$（$m_j\in H^{n_j}(X;\Z/p)$）是这个模的一组基，并考虑映射
\[
f_j:X\longrightarrow K(\Z/p,n_j),\qquad j=1,2,\dots,k,
\]
使 $f_j^*(u_j)=m_j$，其中 $u_j\in H^{n_j}(K(\Z/p,n_j);\Z/p)$ 是复形 $K(\Z/p,n_j)$ 的基本类。于是得到映射
\[
f=\prod_{j=1}^{k}f_j:X\longrightarrow\prod_{j=1}^{k}K(\Z/p,n_j).
\]
把映射 $f$ 化为全空间为 $X_{-1}\sim X$、纤维为 $X_0=F\subset X_{-1}$ 的纤维化：
\[
\tilde f:X_{-1}\longrightarrow\prod_{j=1}^{k}Y_j,\qquad Y_j\sim K(\Z/p,n_j).
\]
然后以 $X_0$ 代 $X_{-1}$ 重复整个程序，如此下去，便得到空间 $X$ 的滤过
\begin{equation}
X\sim X_{-1}\supset X_0\supset X_1\supset X_2\supset\cdots, \tag{7.12}
\end{equation}
其性质是：$\mathcal{A}_p$-模 $M_{j+1}=H^*(X_j,X_{j+1};\Z/p)$ 对 $j\ge0$ 都是自由的。三元组 $(X_{j-1},X_j,X_{j+1})$ 的上同调正合列中的边缘同态
\[
\delta:H^q(X_j,X_{j+1};\Z/p)\longrightarrow H^{q+1}(X_{j-1},X_j;\Z/p)
\]
给出微分 $d_j:M_{j+1}\to M_j$，从而确定下述自由 $\mathcal{A}_p$-模复形：
\begin{equation}
\cdots\longrightarrow M_{j+1}\xrightarrow{d_j}M_j\xrightarrow{d_{j-1}}M_{j-1}\longrightarrow\cdots\xrightarrow{d_0}M_0\xrightarrow{\varepsilon}M_{-1}. \tag{7.13}
\end{equation}
这里 $M_{-1}\cong H^*(X_{-1};\Z/p)$，而 $\varepsilon$ 定义为同态
\[
H^*(X_{-1},X_0;\Z/p)\longrightarrow H^*(X_{-1};\Z/p)
\]
它由对 $(X_{-1},\varnothing)\to(X_{-1},X_0)$ 的包含诱导。

% [原文 p.111]
特别地，在「稳定维数」中有
\[
M_0=H^*(X_{-1},X_0;\Z/p)\cong H^*\Bigl(\prod_{j=1}^{k}K(\Z/p,n_j);\Z/p\Bigr),
\]
因此在代数分解 (7.13) 中，一切 $M_j$（$j\ge0$）都是自由的（且 $\Ker d_j=\Image d_{j+1}$）；于是它代表 $\mathcal{A}_p$-模 $M_{-1}\cong H^*(X;\Z/p)$ 的一个零调分解。复形 (7.13) 称为 $\mathcal{A}_p$-模 $H^*(X;\Z/p)$ 的 \emph{Adams 分解}，而「稳定空间」或谱的图
\begin{equation}
\resizebox{\textwidth}{!}{%
\begin{tikzpicture}[>=Stealth,scale=0.92,every node/.style={inner sep=2pt}]
\node (n0) at (0,0) {$X\sim X_{-1}$};
\node (n1) at (2.9,0) {$X_0$};
\node (n2) at (5.5,0) {$X_1$};
\node (n3) at (7.6,0) {$\cdots$};
\node (n4) at (9.7,0) {$X_j$};
\node (n5) at (12.3,0) {$X_{j+1}$};
\node (n6) at (14.2,0) {$\cdots$};
\node (b1) at (1.45,-1.15) {$X_{-1}/X_0$};
\node (b2) at (4.2,-1.15) {$X_0/X_1$};
\node (b3) at (8.65,-1.15) {$\cdots$};
\node (b4) at (11.0,-1.15) {$X_{j+1}/X_j$};
\node (b5) at (13.4,-1.15) {$\cdots$};
\draw[->] (n1) -- (n0);
\draw[->] (n2) -- (n1);
\draw[->] (n4) -- (n3);
\draw[->] (n5) -- (n4);
\draw[->] (n6) -- (n5);
\draw[->] (n0) -- (b1);
\draw[->] (n1) -- (b1);
\draw[->] (n1) -- (b2);
\draw[->] (n2) -- (b2);
\draw[->] (n2) -- (b3);
\draw[->] (n4) -- (b3);
\draw[->] (n4) -- (b4);
\draw[->] (n5) -- (b4);
\draw[->] (n5) -- (b5);
\draw[->] (n6) -- (b5);
\end{tikzpicture}}
\tag{7.14}
\end{equation}
称为空间 $X$ 的 \emph{Adams 分解}或 Adams 分解 (7.13) 的\emph{几何实现}。从图 (7.14) 取稳定同伦群，就得到 \emph{Adams 谱序列}。我们指出，对 Eilenberg--MacLane 空间 $K(\Z/p,n_j)$ 的乘积，有下述同构：
\[
\Hom_{\mathcal{A}_p}\Bigl(H^*\Bigl(\prod_j K(\Z/p,n_j);\Z/p\Bigr),\Z/p\Bigr)\cong\pi_*^s\Bigl(\prod_j K(\Z/p,n_j)\Bigr),
\]
而这个同构使人们可以用同调代数非常有效地处理 Adams 谱序列。

借助上述 Cartan、Serre、Adams 的方法（往往还结合另一些更特殊的工具与想法），若干拓扑学家完成了球面稳定与非稳定同伦群的深入计算（Serre、Toda 等人，20 世纪 50 年代初至 60 年代中）。到了 60 年代末、进入 70 年代以后，这些技巧又为广义同调论的技巧所补充（见下文 §8、§9）。

球面的稳定同伦群 $\pi_i^s=\pi_{N+i}(S^N)$（$N>i+1$）以及与之相伴的斜对称环 $\pi_*^s$：
\[
\pi_*^s=\sum_{i\ge0}\pi_i^s,\qquad \pi_i^s\pi_j^s=(-1)^{ij}\pi_j^s\pi_i^s,
\]
在光滑流形论中很重要。最初几个球面稳定同伦群由下表给出：

% [Tutor Note: 表中 $i=4,5$ 两格为 $0$，与 $i=12$ 一样都是稳定同伦群的空缺；而 $i=9$ 的 $\Z/2\oplus\Z/2\oplus\Z/2$ 有三个生成元 $\eta^2\sigma,\eta\mu,\nu$，读表时须对照 (7.15) 的关系式才能定出乘法结构。]

% [原文 p.112]
\begin{center}
\footnotesize
\setlength{\tabcolsep}{2.5pt}%
\begin{tabular}{l*{7}{c}}
\toprule
$i=$ & 0 & 1 & 2 & 3 & 4 & 5 & 6\\
\midrule
$\pi_i^s=$ & $\Z$ & $\Z/2$ & $\Z/2$ & $\Z/24$ & $0$ & $0$ & $\Z/2$\\
\midrule
 & $1$ & $\eta$ & $\eta^2$ & $\nu$ & & & $\nu^2$\\
\bottomrule
\end{tabular}
\end{center}

\begin{center}
\footnotesize
\setlength{\tabcolsep}{2.5pt}%
\begin{tabular}{l*{7}{c}}
\toprule
$i=$ & 7 & 8 & 9 & 10 & 11 & 12 & 13\\
\midrule
$\pi_i^s=$ & $\Z/240$ & $\Z/2\oplus\Z/2$ & $\Z/2\oplus\Z/2\oplus\Z/2$ & $\Z/6$ & $\Z/504$ & $0$ & $\Z/4$\\
\midrule
 & $\sigma$ & $\eta\sigma,u$ & $\eta^2\sigma,\eta\mu,\nu$ & $w$ & $x$ & & $y$\\
\bottomrule
\end{tabular}
\end{center}
所标出的生成元之间的乘法关系与加法关系如下：\footnote{回忆在代数 $\Ext_{\mathcal{A}_2}(\Z/2,\Z/2)$ 中，元素 $2$、$\eta$、$\nu$、$\sigma$ 分别由元素 $h_0$、$h_1$、$h_2$、$h_3$ 代表。}
\begin{equation}
\begin{aligned}
2\eta&=0, &&\eta^3=12\nu, &&\eta\nu=0, &&240\sigma=0,\\
\eta^2\sigma&=\nu^3, &&\eta^3\sigma=0, &&\nu^2u=0, &&3w=\eta\nu,\\
\eta^2\nu&=256x, &&y=\nu w, &&6w=0, &&504x=0.
\end{aligned}
\tag{7.15}
\end{equation}
注意元素 $\eta$、$\nu$、$\sigma$、$\eta\sigma$、$\eta^2\sigma$、$x$ 及其倍数都可由球面 $S^n$ 上的标架法丛实现，即都是同态
\[
J:\pi_j(SO)\longrightarrow\pi_j^s
\]
下的像，其定义见第 4 章 §3。（部分稳定群 $\pi_j(SO)$ 列于第 4 章 §2 的表。）非平凡 $p$-分支 $\pi_j^{(p)}\subset\pi_j^s$ 首次出现于 $j=2p-1$；对 $p>2$ 与 $j\le2p(p-1)-1$，它们如下：
\begin{equation}
\begin{aligned}
\pi_j^{(p)}&\cong\Z/p, && j=2k(p-1),\ k=1,\dots,p-1,\\
\pi_j^{(p)}&\cong\Z/p, && j=2p(p-1)-2,\\
\pi_j^{(p)}&\cong\Z/p^2, && j=2p(p-1)-1.
\end{aligned}
\tag{7.16}
\end{equation}
对 $j\ne2p(p-1)-2$，这些 $\pi_j^{(p)}$ 都包含在 $\Image J$ 中，即 $\pi_j^{(p)}\subset J\pi_j(SO)$。结果表明，$J$ 下的像只构成 $\pi_*^s$ 的一小部分；特别地，群 $\pi_{2p(p-1)-2}^{(p)}$ 含有不属于该像的非平凡元素。
% CHUNK ch03h DONE
